\chapter{Concentration on the growing mesh}\label{u:sec:concentration}

In this section, we prove the concentration estimates used in
\cref{u:sec:coarse}. We use the population identities of
\cref{u:sec:proposal-entries} and the product law of the linked rows from
\cref{u:def:laws}. \Cref{u:lem:gram} controls the two blocks of the empirical
matrix. \Cref{u:lem:weighted} controls the weighted errors that enter the
energy expansion in \cref{u:lem:weighted-energy}.

Throughout, $V=V_{\ex{b}}$ and $N=N_{\ex{b}}$. All probability bounds are
under the linked row law, averaged over the disorder. The equality of
laws in \cref{u:sec:proposal-laws} transports them to the source and
replay experiments. Estimates involving only the paths are proved under
the original law of the Brownian paths, which is also their linked marginal.
Stochastic integrals use the spin-free filtration described there.

\section{Errors and dyadic depth}\label{u:sec:gram}

We first define the errors used in both concentration estimates.
For vectors in $\R^n$, we write
$\langle u,v\rangle_n=n^{-1}\sum_{i=1}^n u_iv_i$.
Let $\mathcal T_j$ and $e_j$ be the query vectors of
\cref{u:sec:proposal}, and define
\begin{equation}\label{u:eq:empirical-errors}
\begin{aligned}
 A_{\mathcal T}&=(\langle\mathcal T_i,\mathcal T_j\rangle_n),
 &q&=(\langle\mathcal T_j,x\rangle_n),
 &v&=(\langle e_j,x\rangle_n),
 &H&=(\langle\mathcal T_i,e_j\rangle_n),\\
 E_A&=A_{\mathcal T}-I,
 &E_q&=q-s,
 &E_v&=v-b,
 &E_H&=H-H_0.
\end{aligned}
\end{equation}
Here $v$ is an empirical vector. We use $\|\cdot\|$ for the Euclidean
norm of a vector and the operator norm of a matrix. Let
$D_*=1+\lceil\log_2N\rceil$ be the dyadic depth. We work in the case
$N\ge1$, which holds eventually; when there are no queries, the empty
matrix formulas are interpreted accordingly.

\begin{lemma}[Fixed Gram stability]\label{u:lem:gram}
For $n\ge1$ and $1\le N<n$, let $R=R_V\ge1$ and
$D_*=1+\lceil\log_2N\rceil$. For a linked row
$Z=(\mathcal T,e,x)$, define
$\widehat\Sigma_n=n^{-1}\sum_iZ_iZ_i^{\mathsf T}$ and
$\Sigma=\E ZZ^{\mathsf T}$. Then
\begin{equation}\label{u:eq:fixed-gram}
 \PP\left\{\|E_A\|\vee\|E_H\|>(16D_*^2)^{-1}\right\}
 \le 2\,9^{2N+1}
 \exp\left\{-\frac{n}{2^{24}RD_*^4\log(2+R)}\right\}.
\end{equation}
The blocks $E_A$ and $E_H$ are coordinate compressions of
$\widehat\Sigma_n-\Sigma$.
\end{lemma}

\begin{proof}
We bound the fourth moment of a row projection using its second moment
and its sub-Gaussian tail. We then use a net to control the empirical
covariance matrix, following the empirical-Gram argument of
\cite[Section~4, proof of Proposition~5]{oai-typical-upper}. The
variance-sensitive estimates below keep the dependence on $R$ needed
on the growing mesh.

Let $a=(a_{\mathcal T},a_e,a_x)$ be a deterministic unit vector, and
write $F=a^{\mathsf T}Z$. By the martingale representation, the
normalized $\mathcal T$ projections have sub-Gaussian proxy $R$ times
the squared coefficient norm. The $e$ projections are Gaussian, and
$x$ is a centered sign. By H\"older's inequality with exponents three,
which allows dependence within a row, we have that $F$ has proxy at
most $3R$. By the
population identities in \cref{u:sec:proposal-entries} and
$(u+v+w)^2\le3(u^2+v^2+w^2)$, we also have that $\E F^2\le3$.

Let $Y=F/\sqrt3$ and $Z'=F/\sqrt{3R}$. These have proxies $R$ and $1$,
respectively, and $\E Y^2\le1$. We split $\E(YZ')^2$ at
$(Z')^2=Q=16\log(2+R)$. The lower part is at most $Q$. For the upper
part, we have
\[
 \E Y^4\le16R^2,\qquad
 \E\bigl[(Z')^4\mathbf 1_{\{(Z')^2>Q\}}\bigr]
 \le2e^{-Q/2}(Q^2+4Q+8).
\]
By Cauchy--Schwarz, the upper part is at most
$4\sqrt2R(Q+3)/(R+2)^4<2$. Hence
$\E(YZ')^2\le20\log(2+R)$ and
$\E F^4\le180R\log(2+R)$. Integrating the Gaussian tail, we also have
$\|F^2\|_{\psi_1}\le18R$, where the Orlicz norm is defined below.
This implies that $X=F^2-\E F^2$ has variance at most
$M^2=180R\log(2+R)$ and Orlicz norm at most $S=36R$.

Let $H=4S\log(2+S)$. We split the $l$th moment at $H$. The lower
part is at most $M^2H^{l-2}$, and the upper part is at most
\[
 2H^l(2+S)^{-4}+2l!(2S)^l(2+S)^{-2}.
\]
For $M,S\ge1$, the two upper terms are bounded by $8H^{l-2}$ and
$8l!H^{l-2}$, respectively. Their sum with the lower part is at most
$13l!M^2H^{l-2}$. Expanding the centered exponential series, we find
\[
 \log\E e^{\theta X}
 \le\frac{26M^2\theta^2}{2(1-H|\theta|)}.
\]
Multiplying the moment generating functions of the rows, with
$\theta=t/(26M^2+Ht)$ and both signs, we obtain the two-sided tail
$2e^{-nt^2/[2(26M^2+Ht)]}$ for the empirical mean. The variance
parameter is $4680R\log(2+R)$, and the scale is
\[
 144R\log(2+36R)\le720R\log(2+R).
\]
At tolerance $\varepsilon_0/2$, where
$\varepsilon_0=(16D_*^2)^{-1}\le1$, the rate is at least
$n\varepsilon_0^2/[40320R\log(2+R)]$.

A deterministic $1/4$ net of the unit sphere in dimension $2N+1$ has
size at most $9^{2N+1}$. The operator norm of a symmetric matrix is at
most twice the largest absolute value of its quadratic form on this
net \cite[Corollary~4.2.13 and Exercise~4.4.3(b)]{vershynin-2018}.
Since $40320\cdot256<2^{24}$, we obtain \eqref{u:eq:fixed-gram}
from the union bound over the net. Here we used that the rows have the
product law of \cref{u:def:laws}, and that coordinate inclusions and
projections have norm one.
\end{proof}

\section{Weighted dyadic concentration}\label{u:sec:weighted}

We next prove the weighted estimate. Four elementary estimates will
control products, stochastic integrals with random brackets, weighted
squares, and empirical means. We use them in \cref{u:lem:weighted},
first for each dyadic node and then for the sum over the tree.
The dyadic triangular decomposition is used in
\cite[Section~4, proof of Proposition~5]{oai-typical-upper}. Here the
node estimates retain their durations before they are summed.

We write $\operatorname{SG}(s^2)$ for the centered bound
$\E e^{\theta X}\le e^{s^2\theta^2/2}$ for every real $\theta$, and let
\[
 \|X\|_{\psi_1}
 =\inf\{s>0:\E e^{|X|/s}\le2\}.
\]
The latter is a norm, so its triangle inequality holds also for
dependent random variables; see \cite[Section~2.7.1]{vershynin-2018}.

\begin{lemma}[Products of sub-Gaussian variables]
\label{u:lem:subgaussian-products}
Let $R\ge1$ and $u,s\ge0$. If $\E Y^2\le u^2$, $Y$ is
$\operatorname{SG}(Ru^2)$, and $Z$ is $\operatorname{SG}(s^2)$, then
\begin{equation}\label{u:eq:subgaussian-products}
\begin{aligned}
 \E(YZ)^2&\le20u^2s^2\log(2+R),\\
 \|YZ\|_{\psi_1}&\le10us\sqrt R,\qquad
 \|YZ-\E YZ\|_{\psi_1}\le20us\sqrt R.
\end{aligned}
\end{equation}
The variables $Y$ and $Z$ may be dependent.
\end{lemma}

\begin{proof}
The sub-exponential product estimate is standard
\cite[Lemma~2.7.7]{vershynin-2018}. We additionally retain the second
moment of $Y$ in the variance estimate.
If $us=0$, the assertions are immediate. Otherwise, we normalize to
$u=s=1$. By a union bound, we have that
$\PP(|YZ|>t)\le4e^{-t/(2\sqrt R)}$. Integrating this tail yields
the two Orlicz bounds. To bound the second moment, we split at
$Z^2=Q=16\log(2+R)$. The lower part is at most $Q\E Y^2$. For the
upper part, we use Cauchy--Schwarz, and the Gaussian tails give
\[
 \E Y^4\le16R^2,\qquad
 \E[Z^4\mathbf 1_{\{Z^2>Q\}}]
 \le2e^{-Q/2}(Q^2+4Q+8).
\]
The upper part is at most $4\sqrt2R(Q+3)/(R+2)^4<2$ for $R\ge1$.
In the last inequality, we used $Q+3\le9(R+2)$ and
$R/(R+2)^3\le1/27$. We note that the bound on the second moment uses
both the second moment of $Y$ and the two tails. This proves
\eqref{u:eq:subgaussian-products}.
\end{proof}

\begin{lemma}[A tail bound from the bracket]\label{u:lem:bracket-tail}
Let $\widetilde M$ be a continuous local martingale starting at zero,
and suppose that $\|\langle\widetilde M\rangle_V\|_{\psi_1}\le A$.
Then, for $t,u>0$,
\begin{equation}\label{u:eq:bracket-tail}
\begin{aligned}
 \PP(|\widetilde M_V|>t)
 &\le\PP(\langle\widetilde M\rangle_V>u)+2e^{-t^2/(2u)},\\
 \|\widetilde M_V\|_{\psi_1}&\le5\sqrt{2A}.
\end{aligned}
\end{equation}
\end{lemma}

\begin{proof}
We stop at bracket level $u$ and use the exponential supermartingale
with parameters $t/u$ and $-t/u$, the stopped stochastic-exponential
argument of \cite[Section~3.5]{karatzas-shreve-1991}. On
$\{\langle\widetilde M\rangle_V\le u\}$, we have the first inequality.
For the second, we take $u=t\sqrt{A/2}$ so that the resulting tail is
bounded by $4e^{-t/\sqrt{2A}}$. Integrating it, we find that
$\|\widetilde M_V\|_{\psi_1}\le5\sqrt{2A}$. If $A=0$, the martingale is
zero and the claim is immediate.
\end{proof}

\begin{lemma}[Weighted squares]\label{u:lem:jensen-weights}
Suppose that each $Y_j$ is $\operatorname{SG}(s^2)$, where $s>0$,
and that $p_j\ge0$ and $\sum_jp_j\le1$. Then
\begin{equation}\label{u:eq:jensen-weights}
 \E\exp\left\{\frac{\sum_jp_jY_j^2}{6s^2}\right\}\le2.
\end{equation}
The variables $Y_j$ may be dependent.
\end{lemma}

\begin{proof}
We use the sub-Gaussian-square estimate
\cite[Lemma~2.7.6]{vershynin-2018} and Jensen's inequality, keeping the
constants explicit.
For a single $Y_j$, we integrate its Gaussian tail and bound the
expectation by $1+2\int_1^\infty t^{-3}\dd t=2$. By Jensen's inequality, with the
unused weight assigned to zero, we obtain \eqref{u:eq:jensen-weights}.
\end{proof}

\begin{lemma}[Variance-sensitive Bernstein bound]\label{u:lem:bernstein}
Suppose that $\E X=0$, $\E X^2\le M^2$,
$\PP(|X|>t)\le2e^{-t/S}$, and $M,S\ge1$.
Let $H=4S\log(2+S)$. Then, for integers $l\ge2$,
\begin{equation}\label{u:eq:bernstein-moments}
 \E|X|^l\le13l!M^2H^{l-2}.
\end{equation}
If $X_1,\ldots,X_n$ are independent copies of $X$, then
\begin{equation}\label{u:eq:bernstein}
\begin{aligned}
 \log\E e^{\theta X}
 &\le\frac{26M^2\theta^2}{2(1-H|\theta|)},
 &&|\theta|<H^{-1},\\
 \PP\left(\left|\frac1n\sum_iX_i\right|>\eta\right)
 &\le2\exp\left\{-\frac{n\eta^2}{2(26M^2+H\eta)}\right\},
 &&\eta>0.
\end{aligned}
\end{equation}
\end{lemma}

\begin{proof}
This is the Bernstein inequality in the form of
\cite[Theorem~2.10 and Corollary~2.11]{boucheron-lugosi-massart-2013}; we include the
calculation to record the constants needed here. In passing from the
tail to the factorial moments, we keep the variance.

Splitting the $l$th moment at $H$ and integrating the upper tail, we have
\[
 \E|X|^l
 \le M^2H^{l-2}+2H^l(2+S)^{-4}
       +2l!(2S)^l(2+S)^{-2}
 \le13l!M^2H^{l-2}.
\]
The boundary term is at most $8H^{l-2}$, and the last term is at most
$8l!H^{l-2}$; here we used $H\ge2S$ and
$\log(2+S)\le(2+S)/2$. Expanding the exponential series, which is
absolutely convergent, and using centering, we conclude that the first
inequality in \eqref{u:eq:bernstein} holds. For the empirical mean, we use independence
to multiply the moment generating functions.
Taking $\theta=\eta/(26M^2+H\eta)$, with both
signs, gives the second inequality.
\end{proof}

\begin{lemma}[Duration-normalized weighted concentration]
\label{u:lem:weighted}
Let $m_0=2^{\lceil\log_2N\rceil}$ and $D_*=1+\log_2m_0$.
Pad the errors and $s,b$ by zeros, and $A_{\mathcal T}$ by an identity,
to dimension $m_0$. Dummy coordinates have duration zero. Let
$\mathcal D$ be the complete dyadic tree on $\{1,\ldots,m_0\}$, and
let $P_I$ denote coordinate projection onto $I$. Define
\begin{equation}\label{u:eq:weighted-definitions}
\begin{aligned}
 d_I&=\sum_{j\in I\cap[1,N]}h_{j-1},
 &U_{\rm dur}&=t_N,\\
 \mathcal N_w(X)^2&=\sum_{I\in\mathcal D}\|XP_Iw\|^2,
 &L&=\Upt(E_A),
\end{aligned}
\end{equation}
where $\Upt$ keeps the strict upper triangle and half its diagonal, and
let
\[
\begin{aligned}
 W_{\rm wt}^2={}&\|E_q\|^2+\|E_v\|^2
 +\mathcal N_s(L)^2+\mathcal N_s(L^{\mathsf T})^2\\
 &+\mathcal N_b(L)^2+\mathcal N_s(E_H)^2
 +\mathcal N_s(E_H^{\mathsf T})^2.
\end{aligned}
\]
The duration $d_I$ is that of the previous posterior cells; the query
endpoint $e_j$ has denominator $\sqrt{h_j}$ from the next cell.
Let $R=R_V\ge1$, $B_*=\max(1,\beta e^{1/8})$, and
\begin{equation}\label{u:eq:weighted-parameters}
\begin{aligned}
 \ell_R&=\log(2+R),
 &M_*^2&=25B_*^2D_*^2\ell_R,\\
 S_*&=20B_*D_*\sqrt R,
 &\mathsf v&=26M_*^2,
 &H_*&=4S_*\log(2+S_*).
\end{aligned}
\end{equation}
For every $n,N\ge1$, every deterministic grid with used cells at most
one and $h_j/h_{j-1}\le e^{1/4}$, and every $\eta>0$, under the
linked row product law,
\begin{equation}\label{u:eq:weighted-tail}
 \PP\left\{W_{\rm wt}>2\eta\sqrt{2+5D_*U_{\rm dur}}\right\}
 \le2(10m_0-3)5^{m_0}
 \exp\left\{-\frac{n\eta^2}{2(\mathsf v+H_*\eta)}\right\}.
\end{equation}
With this failure bound, the following estimates hold simultaneously:
\begin{equation}\label{u:eq:weighted-node-bounds}
\begin{aligned}
 &\|E_q\|,\|E_v\|\le2\eta,\\
 &\|LP_Is\|,\|L^{\mathsf T}P_Is\|,\|LP_Ib\|,
   \|E_HP_Is\|,\|E_H^{\mathsf T}P_Is\|
   \le2\eta\sqrt{d_I}\quad(I\in\mathcal D).
\end{aligned}
\end{equation}
Zero-duration nodes have zero error and are omitted when dividing by
$\sqrt{d_I}$. On the growing grid with $h=n^{-\ex{a}}$,
$V=V_{\ex{b}}$, and $\ex{a}>\ex{b}/4$, there exists a fixed finite $C$
such that
\begin{align}
 \PP\{W_{\rm wt}>C_\beta n^{-\ex{w}}\log^C n\}
 &\le2\exp\left\{C_\beta n^{\ex{a}}
 -\frac{c_\beta n^{1-\max\{2\ex{w},\ex{w}+\ex{b}/4\}}}
 {\log^C n}\right\},\label{u:eq:weighted-growing}\\
 \PP\{\|E_A\|\vee\|E_H\|>(16D_*^2)^{-1}\}
 &\le2\exp\left\{C_\beta n^{\ex{a}}
 -\frac{c_\beta n^{1-\ex{b}/2}}{\log^C n}\right\}.
 \label{u:eq:gram-growing}
\end{align}
Constants are independent of deterministic row directions, tree nodes,
the selected cut, and the rotation angle. The finite row estimate allows
all $N$; its transport to the $n$-spin proposal uses $N<n$, which holds
eventually on the parameter family since $\ex{a}<1$.
\end{lemma}

\begin{proof}
We first estimate projections of a single row. We decompose each
triangular projection by dyadic depth and separate its half diagonal.
After dividing each interval projection by $\sqrt{d_I}$, we use
\cref{u:lem:bernstein} and a sphere net. Summing the resulting squared
bounds over the tree gives \eqref{u:eq:weighted-tail}.

\paragraph{Row projections and the durations.}
For a deterministic $\alpha$, the population identities and martingale
representation in \cref{u:sec:proposal-entries} imply that
\begin{equation}\label{u:eq:row-projections}
\begin{aligned}
 \E(\alpha^{\mathsf T}\mathcal T)^2&=\|\alpha\|^2,\\
 \alpha^{\mathsf T}\mathcal T&\text{ is }
       \operatorname{SG}(R\|\alpha\|^2),\qquad
 \alpha^{\mathsf T}e\text{ is }
       \operatorname{SG}(\|\alpha\|^2).
\end{aligned}
\end{equation}
Indeed, $\alpha^{\mathsf T}\mathcal T$ is the integral with integrand
$\alpha_j\chi_t/s_j$ on previous cell $j$; its bracket is bounded by
$R\|\alpha\|^2$. For $u_j=w_js_j$ with $|w_j|\le B_*$, we have
\begin{equation}\label{u:eq:weighted-row-increments}
 Z_J=\sum_{j\in J}u_j\mathcal T_j
 =\sum_{j\in J}w_j\int_{t_{j-1}}^{t_j}\chi_t\dd M_t
 \text{ is }\operatorname{SG}(B_*^2d_I),\qquad J\subset I.
\end{equation}
The bracket bound is deterministic. This implies that the same estimate
holds conditionally before its supporting future cells. Further,
$\|P_Is\|^2\le d_I$, so $s^{\mathsf T}P_Ie$ is
$\operatorname{SG}(d_I)$. The weights $b_j$ satisfy
\begin{equation}\label{u:eq:adjacent-duration-weights}
 b_j\le\beta\sqrt{h_jr(t_j)},\qquad
 s_j\ge\sqrt{h_{j-1}r(t_j)},\qquad
 \frac{b_j}{s_j}\le\beta\sqrt{\frac{h_j}{h_{j-1}}}
 \le\beta e^{1/8}.
\end{equation}
Here we used the monotonicity of $r$ and $p\le1$. Note that the
previous and next durations enter this calculation separately.

\paragraph{The strict triangular part.}
Fix a deterministic unit direction $\alpha$ in the padded dimension
and one node $I$. At a dyadic depth, let $(A,B)$ range over the pairs of
disjoint left and right children of its parents. The projections of
the strict upper triangle at this level and of its transpose are,
respectively,
\[
\begin{aligned}
 N_+&=\sum_{(A,B)}Y_AZ_{B\cap I},
 &Y_A&=\sum_{j\in A}\alpha_j\mathcal T_j,\\
 N_-&=\sum_{(A,B)}Z_{A\cap I}Y_B,
 &Y_B&=\sum_{j\in B}\alpha_j\mathcal T_j.
\end{aligned}
\]
The first is the predictable integral with multiplier $Y_Aw_j\chi_t$
on future cells $j\in B\cap I$. The second has multiplier
$Z_{A\cap I}\alpha_j\chi_t/s_j$ on future cells $j\in B$.
The future supports at one depth are disjoint. By conditional centering
before each future support, the block products are orthogonal, although
their past and future factors need not be independent. The isometry for
$N_+$ and \cref{u:lem:subgaussian-products} for each product in $N_-$
give
\begin{equation}\label{u:eq:dyadic-level-variance}
 \E N_+^2\le B_*^2d_I\sum_A\|P_A\alpha\|^2\le B_*^2d_I,
 \qquad
 \E N_-^2\le20B_*^2d_I\ell_R.
\end{equation}
Their brackets satisfy
\begin{equation}\label{u:eq:dyadic-level-brackets}
 \langle N_+\rangle\le B_*^2d_I\sum_AY_A^2,\qquad
 \langle N_-\rangle
 \le R\sum_B\|P_B\alpha\|^2Z_{A\cap I}^2.
\end{equation}
Using \cref{u:lem:jensen-weights} with weights $\|P_A\alpha\|^2$
and $\|P_B\alpha\|^2$, respectively, we bound the Orlicz norms of
both brackets by $6B_*^2Rd_I$. \Cref{u:lem:bracket-tail} then implies
that the Orlicz norms of both projections are at most
$18B_*\sqrt{Rd_I}$. There are $D_*-1$ depths, and we sum the bounds
over them. By the triangle inequalities in $L^2$ and in the Orlicz
norm, this costs at most a factor equal to the number of depths.

\paragraph{The half diagonal.}
In either orientation, the half diagonal is
\begin{equation}\label{u:eq:half-diagonal}
 Q_I=\frac12\sum_{j\in I}\frac{\alpha_jw_j}{s_j}
       ((\Delta m_j)^2-s_j^2)=N_I+D_I,
\end{equation}
where
\[
\begin{aligned}
 N_I&=\sum_{j\in I}\frac{\alpha_jw_j}{s_j}
 \int_{t_{j-1}}^{t_j}(m_t-m_{t_{j-1}})\chi_t\dd M_t,\\
 D_I&=\frac12\sum_{j\in I}\frac{\alpha_jw_j}{s_j}
 \left(\int_{t_{j-1}}^{t_j}\chi_t^2\dd t-s_j^2\right),
 \qquad \Delta m_j=m_{t_j}-m_{t_{j-1}}.
\end{aligned}
\]
Both terms have mean zero, and $D_I$ is a centered random drift.
By isometry,
\[
 \E N_I^2\le B_*^2\sum_{j\in I}\alpha_j^2h_{j-1}
 \le B_*^2d_I.
\]
Let $S_j=\sup_{[t_{j-1},t_j]}|m_t-m_{t_{j-1}}|$. By the maximal
inequality for the exponential martingale, we have
$\PP(S_j>u)\le2e^{-u^2/(2h_{j-1})}$, so
$\|S_j^2\|_{\psi_1}\le6h_{j-1}$. The bracket of $N_I$ is bounded
by $B_*^2R\sum\alpha_j^2S_j^2$, whose Orlicz norm is at most
$6B_*^2Rd_I$. By \cref{u:lem:bracket-tail}, we have
$\|N_I\|_{\psi_1}\le18B_*\sqrt{Rd_I}$.
For $q_j=\int_{t_{j-1}}^{t_j}\chi_t^2\dd t$, we have
\[
 \|q_j-s_j^2\|_2\le\|q_j\|_2
 \le\int_{t_{j-1}}^{t_j}\sqrt{r(t)}\dd t
 \le\sqrt{h_{j-1}}s_j,
\]
where we used $\chi^4\le\chi^2$. It follows that
\begin{equation}\label{u:eq:half-diagonal-bounds}
\begin{aligned}
 \|D_I\|_2&\le\tfrac12B_*\sqrt{d_I},
 &|D_I|&\le\tfrac12B_*\sqrt{Rd_I},\\
 \E Q_I^2&\le\tfrac94B_*^2d_I,
 &\|Q_I\|_{\psi_1}&\le20B_*\sqrt{Rd_I}.
\end{aligned}
\end{equation}
These estimates keep track of the unequal durations of the cells. In
particular, the variance bound involves $d_I$ without a power of $R$,
which would be lost if each duration were replaced by the maximal one.

\paragraph{The remaining projections and normalization.}
The two projections of the full cross blocks are the centered products
$(\alpha^{\mathsf T}\mathcal T)(s^{\mathsf T}P_Ie)$ and
$(\alpha^{\mathsf T}e)(s^{\mathsf T}P_I\mathcal T)$. By
\eqref{u:eq:row-projections}--\eqref{u:eq:adjacent-duration-weights}
and \cref{u:lem:subgaussian-products}, their variances are at most
$20d_I\ell_R$, and their Orlicz norms are at most $20\sqrt{Rd_I}$.
Here the centering subtracts the block $H_0$.
The two projections without intervals,
$\alpha^{\mathsf T}(\mathcal T x-s)$ and
$\alpha^{\mathsf T}(ex-b)$, have variance at most one and Orlicz
norm at most $10\sqrt R$, since multiplication by the terminal sign
preserves squares and absolute values.

Combining \eqref{u:eq:dyadic-level-variance}--\eqref{u:eq:half-diagonal-bounds},
we find that each of the five interval projections, after division by
$\sqrt{d_I}$, and each of the two projections without intervals is
centered and satisfies
\begin{equation}\label{u:eq:normalized-moments}
 \E X^2\le M_*^2,\qquad \|X\|_{\psi_1}\le S_*.
\end{equation}
For a single row, the five vectors indexed by an interval are
\[
\begin{gathered}
 \Upt(\mathcal T\mathcal T^{\mathsf T}-I)P_Is,\qquad
 \Upt(\mathcal T\mathcal T^{\mathsf T}-I)^{\mathsf T}P_Is,\qquad
 \Upt(\mathcal T\mathcal T^{\mathsf T}-I)P_Ib,\\
 (\mathcal T e^{\mathsf T}-H_0)P_Is,\qquad
 (e\mathcal T^{\mathsf T}-H_0^{\mathsf T})P_Is.
\end{gathered}
\]
To pad these vectors, we extend their errors by zero on the dummy
coordinates. The identity matrix in their row formulas acts on the
first $N$ coordinates.

We have $M_*,S_*\ge1$ in \eqref{u:eq:normalized-moments}, even
for arbitrarily small $d_I$. \Cref{u:lem:bernstein} therefore applies
after the division by $\sqrt{d_I}$, although its hypotheses need not hold
before this normalization.

\paragraph{Nets and the growing grid.}
Fix a deterministic $1/2$ net of the unit sphere of size at most
$5^{m_0}$. For every vector $z$, we have
$\|z\|\le2\max|\alpha^{\mathsf T}z|$, where the maximum is over
this net \cite[Corollary~4.2.13 and Exercise~4.4.2]{vershynin-2018}.
There are $2+5(2m_0-1)=10m_0-3$ vectors to control.
Using \eqref{u:eq:bernstein} and the union bound, we conclude that
\eqref{u:eq:weighted-node-bounds} holds. At every depth, the durations
partition $t_N$, so $\sum_Id_I=D_*U_{\rm dur}$. Summing the squared
errors before normalization, we obtain
\[
 W_{\rm wt}^2\le4\eta^2(2+5D_*U_{\rm dur}),
\]
which proves \eqref{u:eq:weighted-tail}. The union bound is over
$(10m_0-3)5^{m_0}$ events, and the nets are in dimension $m_0$.

For the stopped growing grid, \cref{u:lem:quad} implies that
$N\le8n^{\ex{a}}$, $D_*=O(\log n)$, and
$U_{\rm dur}\le V=O(\log n)$. By \cref{u:lem:tails}, we have
$R\le C_\beta n^{\ex{b}/2}(\log n)^{1/2}$. Then
$M_*^2$ is bounded by a fixed power of $\log n$, and
$H_*\le C_\beta n^{\ex{b}/4}\log^C n$.
Taking $\eta=n^{-\ex{w}}$ in \eqref{u:eq:weighted-tail} and using
$x/(u+v)\ge(1/2)\min(x/u,x/v)$ yields
\eqref{u:eq:weighted-growing}.
To prove \eqref{u:eq:gram-growing}, we use \cref{u:lem:gram} for the
compressions $E_A,E_H$. Its rate is $n/R$ divided by a fixed power of
$\log n$, which gives the separate loss $\ex{b}/2$.
By the equality of row laws in \cref{u:sec:proposal-laws}, all these
bounds also hold under the linked source and replay laws.
\end{proof}
